\section*{\centering ВИСНОВКИ}\addcontentsline{toc}{section}{ВИСНОВКИ}

У дисертаційному дослідженні теоретично обґрунтовано та розроблено модель формування інформаційно-комунікативної компетентності в процесі фахової підготовки майбутніх дизайнерів на основі комплексного обчислювального аналізу 122 українських освітніх програм спеціальності 022~Дизайн (3\,043 освітні компоненти), поглибленого кейс-стаді трьох програм КДПУ (132~силабуси, опитування 32~випускників, 9~роботодавців, 61~студента) та міжнародного бенчмаркінгу 7~провідних програм із 6~країн. Відповідно до поставлених завдань дослідження отримано такі основні результати.

\begin{enumerate}

\item \textit{Визначено теоретичні засади формування ІКК у контексті трансформації дизайн-освіти ери ШІ} (розділ~\ref{sec:1}).

Термінологічний аналіз наукових джерел засвідчив, що інформаційно-комунікативна компетентність є найбільш адекватним поняттям для описання інформаційно-комунікативної готовності дизайнерів, оскільки інтегрує як інформаційні, так і комунікативні аспекти професійної діяльності, на відміну від суміжних понять (ІКТ-компетентність, цифрова компетентність, медіакомпетентність). На цій основі обґрунтовано 4-компонентну структуру ІКК, специфічну для дизайнерів: інформаційно-аналітичний (пошук, аналіз, оцінювання інформації, у~т.\,ч. згенерованої ШІ), комунікативний (цифрова комунікація, включаючи промпт-інженерію як нову форму людино-машинної взаємодії), технологічний (застосування цифрових та ШІ-інструментів дизайну, включаючи генеративний ШІ), рефлексивний (самооцінка власного рівня ІКК та етичні аспекти використання технологій). Оглядове дослідження (scoping review) 156~наукових праць з проблеми генеративного ШІ в дизайн-освіті (Web of Science, 2017--2026) виявило 6,4-кратне зростання публікацій після появи ChatGPT (86,5\% досліджень опубліковано після 2022~р.), зафіксувало парадигмальний зсув від <<дизайнера-виробника>> до <<дизайнера-куратора>> та встановило критичну теоретичну прогалину~--- лише 26,3\% досліджень використовували теоретичні рамки (73,7\% атеоретичних). Зіставлення 4-компонентної моделі ІКК із рамкою DigComp~2.2 та стандартом вищої освіти В2 спеціальності 022 у контексті Європейської рамки кваліфікацій підтвердило структурну відповідність за вимірами знань, умінь та автономності, проте виявило суттєві прогалини у компетентностях, пов'язаних із критичною рефлексією, міждисциплінарною співпрацею, сталим дизайном та цифровими/ШІ-навичками. Розроблено інтегративну теоретичну рамку дослідження, що поєднує три підходи: модель TPACK (адаптовану для дизайну: TK\,=\,ШІ/цифрові інструменти, PK\,=\,студійні методи, CK\,=\,предметна область дизайну), студійну педагогіку (рефлексивна практика Шона) та компетентнісний підхід (стандарт В2). Ця рамка є наскрізною для всього дослідження: вона спрямовує обчислювальний аналіз даних та обґрунтовує модель трансформації навчальних програм.

\item \textit{Проведено комплексний аналіз поточного стану формування ІКК в українській дизайн-освіті на основі даних} (розділ~\ref{sec:2}).

Загальнонаціональний обчислювальний аналіз 122 акредитованих програм спеціальності 022~Дизайн (3\,043 освітні компоненти, 1\,515 унікальних назв дисциплін) виявив системну прогалину у формуванні ІКК: при задовільному покритті традиційних цифрових інструментів (71,4\%) компоненти <<Основи ШІ>>, <<Генеративний ШІ>> та <<Комп'ютерний зір>> мають 0\% покриття, а загальний рівень ШІ-розриву становить 35,7\%. Тематичне моделювання (BERTopic) засвідчило, що домінуючі теми зосереджені навколо традиційного мистецтва, практичної підготовки та загальноосвітніх дисциплін~--- жодна тема не стосується ШІ чи цифрових інновацій. Мережевий аналіз навчальних планів (1\,240 вузлів, 7\,524 ребер, 10~спільнот Лувена) підтвердив, що цифрові та ШІ-дисципліни не утворюють самостійного кластера, розчиняючись у периферії інших спільнот. Кластеризація за семантичною подібністю виявила 5~програмних кластерів із високою внутрішньою однорідністю та традиційною орієнтацією, що свідчить про системну ригідність навчальних планів. Кейс-стаді трьох програм КДПУ (графічний дизайн, дизайн одягу, дизайн середовища) із тріангуляцією 6~джерел даних виявив нульове покриття ШІ-тематики у 132~проаналізованих силабусах, нисхідний тренд ІКК-ключових слів в ОПП ($-$35\% між 2017 та 2021~рр.), незадоволеність випускників ІКТ-підготовкою (33,3\% у графічному дизайні) та пряму рекомендацію роботодавців щодо інтеграції ШІ. Статистичний аналіз за критерієм Краскела--Уолліса виявив значущі міжпрограмні відмінності у задоволеності матеріально-технічним ($H = 8{,}708$, $p = 0{,}013$) та інформаційно-методичним ($H = 6{,}019$, $p = 0{,}049$) забезпеченням. Рефлексивний компонент визначено як найбільшу прогалину у формуванні ІКК (11--12 згадувань експертами проти 70--114 для інших компонентів). Міжнародний бенчмаркінг 7~провідних програм дизайну (Великобританія, США, Фінляндія, Нідерланди, Австралія, Італія, Південна Корея) підтвердив суттєвий розрив: міжнародні програми містять від 1 до 4~ШІ-дисциплін та інтегрують промпт-інженерію, ШІ-етику і критичний дизайн, тоді як жодна з 122~українських програм не має таких дисциплін.

\item \textit{Розроблено та обґрунтовано модель трансформації навчальних програм для формування ІКК} (розділ~\ref{sec:3}).

На основі інтегративної теоретичної рамки (розділ~\ref{sec:1}) та емпіричних даних (розділ~\ref{sec:2}) розроблено структурно-функціональну модель формування ІКК майбутніх дизайнерів, що складається з 5~взаємопов'язаних блоків: мотиваційно-цільового (соціальне замовлення, зумовлене 0\% покриттям ШІ у 122~програмах, незадоволеністю випускників до 33,3\%, рекомендацією роботодавців щодо ШІ), теоретико-методологічного (TPACK + студійна педагогіка + компетентнісний підхід; 6~дидактичних принципів), змістового (4~компоненти ІКК, відображені на навчальні дисципліни), організаційно-технологічного (8~категорій ШІ-інструментів, модифіковані студійні методи) та діагностично-результативного (4~критерії, портфоліо-оцінювання, 4~рівні сформованості ІКК). Розроблено інтегровану модель 240-кредитного навчального плану (4~роки, 8~семестрів) із наскрізною ШІ-інтеграцією, що передбачає 4~нові ШІ-дисципліни (<<Цифрова грамотність та основи AI>>, <<Генеративний дизайн та AI>>, <<Взаємодія людина-AI>>, <<Критичний дизайн та етика AI>>), при цьому 76,7\% дисциплін (33~з~43) можуть інтегрувати ШІ на базовому, інструментальному або трансформаційному рівнях із збереженням фундаментальних дисциплін (Рисунок, Живопис, Композиція). Обґрунтовано 4~педагогічні умови формування ІКК: системна інтеграція ШІ наскрізно через навчальний план (відповідь на 0\% покриття ШІ та принцип TPACK), трансформація студійної педагогіки від <<виробника>> до <<виробника + куратора>> (відповідь на парадигмальний зсув, зафіксований в оглядовому дослідженні), узгодження з вимогами роботодавців (відповідь на пряму рекомендацію щодо ШІ та незадоволеність випускників), безперервний моніторинг ІКК через портфоліо-оцінювання (відповідь на недостатність рефлексивного компоненту). Розроблено 4-етапну стратегію поетапного впровадження моделі протягом 4~років навчання: усвідомлення та підготовка (семінари для викладачів, пілотна інтеграція), інструментальна інтеграція (оновлення 8--10 силабусів, модифікація критики), поглиблена інтеграція (2~нові ШІ-дисципліни, ШІ-лабораторія, трансформаційний рівень), повна реалізація та оцінка (<<Критичний дизайн та етика AI>>, ШІ у дипломному проєкті, комплексна оцінка).

\end{enumerate}

Таким чином, гіпотезу дослідження підтверджено: формування ІКК майбутніх дизайнерів може бути системно вдосконалено на основі визначення теоретичних засад ІКК у контексті ШІ-трансформації дизайн-освіти, комплексного обчислювального аналізу програм із застосуванням тематичного моделювання, мережевого аналізу та аналізу розривів у загальнонаціональному масштабі (122~програми, 3\,043~освітні компоненти), а також розробки моделі трансформації навчальних програм на основі даних, що визначає інтеграцію ШІ-компетентностей для кожного типу навчальних дисциплін.

Перспективними напрямами подальших досліджень є: експериментальна верифікація запропонованої моделі формування ІКК через формувальний педагогічний експеримент; розробка діагностичного інструментарію для вимірювання рівня сформованості ІКК майбутніх дизайнерів за 4~компонентами та 4~рівнями; лонгітюдне дослідження ефектів ШІ-інтеграції на якість професійної підготовки та конкурентоспроможність випускників на ринку праці; поширення методики обчислювального аналізу та моделі трансформації навчальних програм на інші творчі спеціальності (архітектура, образотворче мистецтво, аудіовізуальне мистецтво); міжнародне порівняльне дослідження моделей ШІ-інтеграції в дизайн-освіті.

\label{end_main_text}
